\section{CQ-to-RA administrative algorithms}
\label{sec:planning}
The graph and prefix algorithms are distinct transition systems inside the
same existing pass. A graph-domain judgment selects which is applicable; it
is not an alpha-acyclicity test and not an additional IR. This section specifies
both algorithms, their deterministic policies, and their semantic domains.

\subsection{Dependency and movement judgments}
For clause $c_i$, derive
$\Sigma;B;T\vdash c_i\triangleright(R_i,O_i,m_i,a_i,K_i)$:
required names, provided names, movable bit, optional atom arity, and literal
positions. $B$ is the written prefix and $T\subseteq B$ contains names bound
by preceding lets. $O_i$ includes already-bound positive-atom equality keys;
it is not always a set of freshly introduced names.
\begin{mathpar}

\Rule{D-Guard}{R=\uses(e,B)}
 {\Sigma;B;T\vdash\mathsf{if}\ e\triangleright(R,\emptyset,m_e,\bot,\emptyset)}

\Rule{D-Bind}{N=\set(\binds(p))\quad R=\uses(e,B)}
 {\Sigma;B;T\vdash\mathsf{let}\ p[:\tau]=e\triangleright(R,N,m_{p,e},\bot,\emptyset)}
\end{mathpar}
For D-Atom, walk arguments $j$ left to right with initially empty $L,R,K$.
A bare $x$ adds $x$ to $L$. A literal adds $j$ to $K$. Wildcard adds nothing.
A computation $e$ adds $\uses(e,B\cup L)\setminus L$ to $R$ and makes the
atom a barrier. Conclude
\[
 \Rule{D-Atom}{\operatorname{walk}(\bar a,B)=(L,R,K,m)}
 {\Sigma;B;T\vdash R_0(\bar a)\triangleright
 (R\cup(L\cap T),L,m,|\bar a|,K)}.
\]
Update $B$ by $O_i$ after each clause, and $T$ by the names of a let.
Thus a preceding let is a required value producer even when an atom later
uses the same bare name. An earlier positive atom is not a unique producer
of a shared logical join variable.

D-Move assigns $m$ as follows. Positive atoms are movable exactly when every
argument is bare, wildcard, or literal. A let is movable only for a single
identifier pattern with neither by-reference nor subpattern modifier and a
literal RHS. A guard is movable only if it is recognized as one of six
variable/integer comparisons and every operand variable has syntactically
primitive i32 type at \emph{every} occurrence/binding. All other clauses have
$m=\mathsf{false}$. Type accumulation precedes final guard classification;
Rust aliases are not resolved by this analysis. These are current algorithmic
premises, not a general effect inference system.

D-Graph holds exactly when the head is distinct bare names supplied by the
body, every positive atom has distinct bare arguments, there is no let, and
all guards pass D-Move. This admits cyclic and disconnected hypergraphs.
\begin{mathpar}
\Rule{D-Move}{\operatorname{simpleAtom}(c)\ \lor\operatorname{literalBind}(c)\ \lor\operatorname{primitiveGuard}(c)}
 {\Sigma\vdash c:\mathsf{movable}}

\Rule{D-Barrier}{\neg(\operatorname{simpleAtom}(c)\lor\operatorname{literalBind}(c)\lor\operatorname{primitiveGuard}(c))}
 {\Sigma\vdash c:\mathsf{barrier}}

\Rule{D-Graph}{\operatorname{distinctBareHead}(q)\quad\operatorname{headBound}(q)\\
 \forall c\in q.\ c\text{ is a distinct-bare atom or primitive guard}}
 {\Sigma\vdash q:\mathsf{graphDomain}}
\end{mathpar}
The three movement predicates are exactly the cases defined above; they are
not inferred from the absence of visible free names.
Negation and aggregation have no D-rule in the implemented semantic domain.

\subsection{SIP graphs, query hypergraphs, and the GYO variant}
For $n\ge1$, let $A_0,\ldots,A_{n-1}$ be relational body occurrences and
$X_i=\vars(A_i)$. The \emph{sideways information-passing (SIP) graph}
has an edge $\{i,j\}$ exactly when $i\ne j$ and $X_i\cap X_j\ne\emptyset$;
constant-containing atoms are marked. Constraint atoms contribute no vertices
(AHV Section~6.1, p.~113; note 2.2). The compiler does not construct a separate
SIP-graph IR. Clause producer requirements and movement barriers in D-* are
additional compiler facts, not the definition of a SIP graph.

AHV's hypergraph has distinct nonempty hyperedges (Section~6.4, p.~130).
Note 2.3's \emph{query hypergraph} keeps an occurrence label for each $X_i$.
We use that course adaptation: equal sets retain distinct occurrence identities.
Identifying repetitions and discarding empty sets gives the underlying
textbook hypergraph $(\bigcup_iX_i,\{X_i\mid X_i\ne\emptyset\})$.
Empty sets in GYO's residual state are algorithmic objects, not extra edges
in AHV's definition. Empty original sets, admitted by the helper's general
domain, also keep their separate occurrence identities.

A \emph{join tree} $J$ is a tree on occurrences with each edge $\{i,j\}$
labeled by $X_i\cap X_j$. Every variable shared by two nodes must appear on
every edge of their path. Equivalently, its containing nodes induce a connected
subtree. This is AHV's join-tree condition (p.~129) applied to the occurrence
family. By Theorem~6.4.5 it characterizes AHV acyclicity; we use the notes'
term \emph{alpha-acyclic} for this notion. Such a tree is different from a
binary logical join plan whose internal nodes are joins of intermediate results.

AHV Algorithm~6.4.4 removes \emph{ears}: an edge $f$ with witness $g\ne f$
satisfies $f\cap\bigcup(F\setminus\{f\})\subseteq g$; an isolated edge is
also an ear. It succeeds when all edges have been removed. The following
private-variable/absorption system is the compiler's GYO variant, with its
correspondence to ear removal proved below.

GYO state is $(Z,\pi)$, with a map of active identities to residual variable sets
and a parent vector for absorbed identities. Initially $Z(i)=X_i$ and
$\pi(i)=\bot$. Neither $X$ nor any executable occurrence changes.
\begin{mathpar}

\Rule{G-Private}{i\in\dom Z\quad x\in Z(i)\quad
 \forall j\in\dom Z\setminus\{i\}.\ x\notin Z(j)}
 {(Z,\pi)\step{G-Private}(Z[i\mapsto Z(i)\setminus\{x\}],\pi)}

\Rule{G-Absorb}{i,j\in\dom Z\quad i\ne j\quad Z(i)\subseteq Z(j)}
 {(Z,\pi)\step{G-Absorb}(Z\setminus\{i\},\pi[i\mapsto j])}
\end{mathpar}
Policy selects the lexicographically first applicable Private pair $(i,x)$,
otherwise the first Absorb pair $(i,j)$. A terminal state succeeds exactly
when there is one active occurrence with an empty residual set; its identity is the sole root and
$\pi$ supplies the tree. All other terminal states report no tree.
Writing $\operatorname{terminal}(Z)$ for absence of either GYO step:
\begin{mathpar}
\Rule{G-Accept}{\operatorname{terminal}(Z)\quad Z=\{i\mapsto\emptyset\}}
 {(Z,\pi)\Downarrow_{\rm GYO}\mathsf{Some}(\pi)}

\Rule{G-Reject}{\operatorname{terminal}(Z)\quad Z\ne\{i\mapsto\emptyset\}\text{ for every }i}
 {(Z,\pi)\Downarrow_{\rm GYO}\mathsf{None}}
\end{mathpar}
The certificate is replayed, and the original variable sets independently undergo
connectivity and running-intersection checking.

\begin{lemma}[Correspondence to AHV ear removal]
For distinct nonempty original edges, every absorption is a legal AHV ear
removal on the remaining original edges. A successful compiler derivation,
followed by removal of its sole remaining original edge, gives a successful
AHV derivation.
\end{lemma}
\begin{proof}
If two original edges containing $x$ remain active, both residuals still
contain $x$: otherwise consider its first deletion from either residual.
Private could not have deleted it while it remained in the other residual.
Absorb only removes whole occurrences. Thus for distinct active $i,j$,
$X_i\cap X_j=Z(i)\cap Z(j)$. At absorption of $i$ into $j$, every variable
of $X_i$ shared with another active edge lies in $Z(i)\subseteq Z(j)\subseteq X_j$.
This is exactly the ear-witness condition on original edges. Private steps
leave those original edges unchanged. The last edge is isolated and hence an
ear. Equal-edge occurrence deletions and empty original edges stutter when
projected to the textbook hypergraph; the next theorem proves recognition
and the certificate directly on the full occurrence domain.
\end{proof}

\begin{theorem}[GYO termination and recognition]
Every GYO derivation terminates. It succeeds iff $X$ admits a join tree.
On success its absorption parent vector is a join tree for the original
occurrences.
\end{theorem}
\begin{proof}
The natural-number measure $|\dom Z|+\sum_i|Z(i)|$ strictly decreases on each
step. Private deletion preserves existence of a join tree: removal cannot
break connectedness of any other variable's nodes; restoration adds a
variable to only one active node. Absorption also preserves existence. To
restore a contained edge $i\subseteq j$, attach $i$ as a leaf to $j$. To delete
$i$ from an existing join tree, remove its node, retain the component containing
$j$, and attach every other resulting component to $j$ at its former neighbor
of $i$. A variable shared across such a cut occurs in $i$, hence in $j$;
its path inside each component contains its attachment node. Thus running
intersection is preserved, and the new graph is a tree.

If a join tree has at least two nodes, choose a leaf $i$ and its neighbor $j$.
Any variable of $X_i\setminus X_j$ occurs nowhere else by running intersection,
so a Private rule exists unless that difference is empty. If it is empty,
Absorb applies. A one-node nonempty hypergraph has a Private rule. Therefore
an acyclic terminal state is one empty edge. Conversely that state admits a
join tree, and reversing all steps reconstructs one for the original family.
Every reversed Absorb attaches the recorded child to its recorded parent;
reversed Private changes no tree edge. The reconstruction is exactly $\pi$.
The argument includes duplicate and empty edges and trees joined by empty
separators.
\end{proof}

\subsection{Assembly, selection, and retention}
Let $\bar h$ be the ordered head variables, $H=\set(\bar h)$, and $P$ the
set of pending predicate identities. Variable names denote their corresponding
attributes as fixed in Section~\ref{sec:valuations}.
An assembly frontier maps group identifiers to $(C,V,r)$, where $C$ is
occurrence coverage, $V$ the sort, and $r$ the current immutable result.
Coverages are nonempty, pairwise disjoint, and partition all introduced
occurrences. For complete introduction define
\[
 K(C,P)=X(C)\cap\left(H\cup X(\bar C)\cup\vars(P)\right),\qquad
 X(C)=\bigcup_{i\in C}X_i.
\]
Here $K$ is note 2.4's retained-sort function \code{keep}; the notes' pending-use
set is $\vars(P)$ in this state's notation.
The invariant requires $K(C,P)\subseteq V\subseteq X(C)$ and
\begin{equation}
\label{eq:global}
 \operatorname{out}_{\bar h}\!\left(
 \pi_H\sigma_{\bigwedge P}\left(\mathop{\bowtie}_{(C,V,r)\in F}\sem r\right)\right)
 =\sem q_I.
\end{equation}
Predicates already applied are reflected in $\sem r$; they must not be omitted
from an alleged per-group denotation. A proof may track their group ownership
as ghost information, without adding it to the executable IR.

$\emit(D,o)=(D\concat[r=o],r)$ chooses a fresh result and checks its resulting
sort against any configured arity cap before updating state. P-Start emits
a full positional Rename for the next occurrence $i$, inserts
$(\{i\},X_i,r_i)$, and advances the introduction counter. After all $n$
introductions the phase is Local. The introduction transition, with counter $i$, is
\[
\Rule{P-Start}{\phi=\mathsf{Starting}\quad 0\le i<n\quad
 \emit(D,\delta_{m_i}(R_i))=(D',r_i)\quad g_i\text{ fresh in }F}
 {(D,F,P,i,\phi)\step{P-Start}
 (D',F[g_i\mapsto(\{i\},X_i,r_i)],P,i+1,\phi')},
\]
where $g_i=i$ and $\phi'=\mathsf{Local}$ exactly when $i+1=n$,
otherwise Starting. The remaining rules are:
\begin{mathpar}

\Rule{P-Select}{\phi\in\{\mathsf{Local},\mathsf{Assembly}\}\quad
 p\in P\quad F(g)=(C,V,r)\quad\vars(p)\subseteq V\quad
 \emit(D,\sigma_pr)=(D',s)}
 {(D,F,P,\phi)\step{P-Select}(D',F[g\mapsto(C,V,s)],P\setminus\{p\},\phi)}

\Rule{P-Begin}{\phi=\mathsf{Local}\quad\text{all occurrences introduced}\quad
 \nexists\text{ applicable P-Select}}
 {(D,F,P,\phi)\step{P-Begin}(D,F,P,\phi_J)}

\Rule{P-Combine}{F(g)=(C,V,r)\quad F(h)=(E,W,s)\quad C\cap E=\emptyset\\
 \phi=\mathsf{Assembly}\quad\text{no ready Select or Trim}\quad
 \emit(D,r\bowtie s)=(D',t)}
 {(D,F,P,\phi)\step{P-Combine}
 (D',(F\setminus\{g,h\})[g\mapsto(C\cup E,V\cup W,t)],P,\phi)}

\Rule{P-Trim}{F(g)=(C,V,r)\quad U=K(C,P)\subsetneq V\quad
 \phi=\mathsf{Assembly}\\\text{no ready Select}\quad\emit(D,\pi_Ur)=(D',s)}
 {(D,F,P,\phi)\step{P-Trim}(D',F[g\mapsto(C,U,s)],P,\phi)}

\Rule{P-Output}{\phi=\mathsf{Assembly}\quad F=\{(\{0,\ldots,n-1\},H,r)\}
 \quad P=\emptyset}
 {(D,F,P,\phi)\step{P-Output}(D,\mathsf{output}\ r\ \mathsf{as}\ h(\bar x))}
\end{mathpar}
$\phi_J$ is Reduction when a checked join tree is selected and Assembly
otherwise. Reduction initialization copies frontier references without
projecting their sorts. P-Trim specifies the implemented eager policy;
the more general safe rule permits any $K(C,P)\subseteq U\subsetneq V$.
Output's printed group map omits its arbitrary sole group identifier.

\begin{lemma}[Retention]
Suppose $L:V$, $M:W$, and $V\cap(W\cup H\cup\vars(P))\subseteq U\subseteq V$.
Then
$\pi_H\sigma_P(L\bowtie M)=\pi_H\sigma_P(\pi_UL\bowtie M)$
whenever the displayed head and predicates are well-scoped.
\end{lemma}
\begin{proof}
Left-to-right, restrict the left witness; every common coordinate and every
observed/predicate coordinate is retained, so compatibility and tests persist.
Right-to-left, a tuple of $\pi_UL$ has a witness in $L$. Extend it by that
witness. All coordinates shared with $M$ are in $U$, so the witness remains
compatible. Its head and predicate coordinates are unchanged. Both sides
therefore contain exactly the same head tuples, including the empty-head case.
\end{proof}

\begin{theorem}[Assembly preservation and termination]
On the graph domain under $\mathcal H_{\rm rel}$, every successful complete
assembly derivation denotes $q$. The implemented schedule terminates or
returns an arity error.
\end{theorem}
\begin{proof}
Full positional Rename introduction establishes $\Val_I(A_i)$ for each atom;
Eq.~\eqref{eq:global} initially follows from the CQ definition. Select moves
one conjunct into a relation containing all its free names; distributing this
conjunct over the surrounding join gives the same global relation. Combine
uses associativity of compatible tuple union on disjoint coverages. Trim uses
the retention lemma: remaining groups use only outside-occurrence variables,
and all remaining predicates/head variables are in $K$. Boundary variables
shared by any future combine are therefore present on both sides. Output has
one complete coverage and no pending predicate, so Eq.~\eqref{eq:global}
is its exact result. Select decreases $|P|$, Combine decreases $|F|$, and Trim
strictly decreases a finite sort between combinations. The lexicographic
measure $(|F|,|P|,\sum_{g\in F}|V_g|)$ decreases in assembly. Introduction and
the reduction schedule are finite. An arity error is a controlled rejection,
not a proof that no plan meeting the cap exists.
\end{proof}

\subsection{Semijoin programs, consistency, and immutable versions}
For occurrence relations $L_i$ of sorts $X_i$, AHV's \emph{pairwise consistency}
means $L_i\ltimes L_j=L_i$ for every $i,j$; \emph{global consistency} means
$L_i=\pi_{X_i}(\bowtie_jL_j)$ for every $i$. A \emph{semijoin program} is a
sequence of semijoin replacements. It is a \emph{full reducer} if it produces
global consistency on every input instance (Section~6.4, pp.~128--129).
Note 2.5's ``complete support'' describes that global-consistency condition.
The immutable definitions below implement those replacements by new result
versions. Pending multi-occurrence guards are not part of the reduced family.

Let $J$ be rooted, $\operatorname{pa}(c)$ its parent, and
$\operatorname{ch}(c)$ its children. Reduction state is $(D,\mathbf r,U,Dn)$,
where $\mathbf r$ maps occurrences to their current RA versions; $U,Dn$ record
completed upward/downward edges by child identity. Let $E_J$ be the nonroot
nodes. Initial $\mathbf r$ refers to the locally selected occurrences $L_i$.
\begin{mathpar}

\Rule{Y-Up}{c\in E_J\setminus U\quad\operatorname{ch}(c)\subseteq U\quad Dn=\emptyset\\
 p=\operatorname{pa}(c)\quad\emit(D,\mathbf r(p)\ltimes\mathbf r(c))=(D',s)}
 {(D,\mathbf r,U,Dn)\step{Y-Up}(D',\mathbf r[p\mapsto s],U\cup\{c\},Dn)}

\Rule{Y-Down}{U=E_J\quad c\in E_J\setminus Dn\quad p=\operatorname{pa}(c)\\
 (p=\operatorname{root}\lor p\in Dn)\quad
 \emit(D,\mathbf r(c)\ltimes\mathbf r(p))=(D',s)}
 {(D,\mathbf r,U,Dn)\step{Y-Down}(D',\mathbf r[c\mapsto s],U,Dn\cup\{c\})}

\Rule{Y-Freeze}{U=Dn=E_J}
 {(D,\mathbf r,U,Dn)\step{Y-Freeze}
 (D,\{(\{i\},X_i,\mathbf r(i))\}_{i<n},\mathsf{Assembly})}
\end{mathpar}
Freeze transfers references; it does not recompute, mutate, or delete earlier
results. The implementation uses reverse preorder for Up and preorder for
Down. Any sibling ordering satisfying these premises is legal.

\begin{theorem}[Full reducer]
After Freeze, with $J$ satisfying running intersection,
$\sem{\mathbf r(i)}=\pi_{X_i}(\bowtie_j L_j)$ for every $i$.
Every reduction step preserves the full join of the current occurrence family.
\end{theorem}
\begin{proof}
A tuple participating in the full join has a compatible support tuple in every
other occurrence. Replacing any factor $A$ by $A\ltimes B$ therefore removes
no full-join witness; the replacement is a subset, so it adds none.

For completeness of support, induct upward on subtrees. After processing all
children of $i$, a tuple of the current $i$ survives exactly when it extends
to a solution in each child subtree. Those extensions can be combined: any
variable shared by two different child subtrees occurs at $i$ by running
intersection, so they agree through the tuple at $i$. Hence the current
relation is the projection of the subtree join onto $X_i$. At the root this
is projection of the complete join. Induct downward. The parent relation has
complete support by the induction hypothesis. A child tuple surviving its
upward reduction has subtree support; its matching retained parent supplies
outside support. Variables common to the child's subtree and its complement
occur at both endpoints of the separator, again by running intersection.
The two supports thus combine. Conversely every full-join child tuple matches
its parent and survives. This proves exact support at the child. Induction
covers all nodes, including empty separators: an empty support relation
propagates emptiness, whereas a nonempty one supplies the unique empty key.
\end{proof}
Pending cross-occurrence guards remain for assembly. The full-support theorem
concerns the locally selected $L_i$, not those unapplied theta predicates.
It makes no projected-output-linear or WCOJ complexity claim.

\subsection{Policies and the canonical route}
Written combines the two smallest frontier identifiers. Connected keeps the
smallest left identifier and chooses the first intersecting right sort,
otherwise the next identifier. GreedyWidth minimizes
$(|V_g\cup V_h|,g,h)$ among $g<h$. Auto first performs GYO; if it finds a tree,
it uses Up, Down, Freeze and grows the root group by an adjacent remaining
occurrence whose parent is already covered. Otherwise Auto uses GreedyWidth.
Available selections precede trims, and both precede the next combination.
These policy functions never emit RA. Following note 2.4, the measured
quantity is \emph{logical arity}, $|\sort(L)|$. Relation cardinality $|L|$,
physical tuple width in bytes, and hypergraph-decomposition widths are separate
quantities. The implementation name \code{GreedyWidth} denotes the stated
arity heuristic; no cardinality estimator or decomposition-width optimizer
is implemented.

The canonical R1 algorithm emits the first full Rename as prefix, then for
each subsequent occurrence emits Rename followed by natural join with the
prefix, and ends with Project to the head. These are P-Start/Combine/Output's
relational operations with a fixed left-deep schedule and only final pruning;
R1 is not required to satisfy the planned kernel's eager-Trim premise.
For source occurrence Rename $a_i=\delta_{m_i}(R_i)$, the exact canonical rules are
\begin{mathpar}
\Rule{B-Root}{p=\bot\quad\emit(D,a_i)=(D',r)}
 {(D,p;i)\Rightarrow_B(D',r;i+1)}

\Rule{B-Join}{p\ne\bot\quad\emit(D,a_i)=(D_1,r)\quad\emit(D_1,p\bowtie r)=(D',s)}
 {(D,p;i)\Rightarrow_B(D',s;i+1)}

\Rule{B-Output}{i=n\quad\emit(D,\pi_Hp)=(D',r)}
 {(D,p;i)\Rightarrow_B(D',\mathsf{output}\ r\ \mathsf{as}\ h(\bar x))}
\end{mathpar}
Its correctness follows directly by induction on the written atom list, then
projection. R1 delegates extended CQ to the prefix rules with Written policy.

\subsection{Prefix transitions for Rust clauses}
State is $S=(D,p,B,U,K,o)$: immutable definitions, optional prefix result,
complete prefix sort, pending clause identities, normalized literal keys,
and optional output. Fresh-name supply is implicit. Initially $D,B,K$ are
empty, $U=\{0,\ldots,n-1\}$, and $p,o$ absent. When an operator first needs a
prefix, $\operatorname{unitize}$ emits $r=\mathbf1$ and returns $r$; otherwise
it returns $p$ without emission. No rule is allowed once $o$ is present.
\begin{mathpar}

\Rule{C-Ready}{\operatorname{normalized}(K)\quad o=\bot\quad i\in U\quad R_i\subseteq B\\
 \forall j\in U.\ j<i\Longrightarrow(m_i\land m_j)\quad
 (c_i\text{ is Bind}\Longrightarrow O_i\cap B=\emptyset)}
 {S\vdash i:\mathsf{ready}}

\Rule{C-Constant}{(i,j)\text{ is the next unnormalized literal position}\\
 U=\{0,\ldots,n-1\}\quad k\text{ fresh}\quad
 (D',p')=\operatorname{typedBind}(D,p,k,\Sigma(R_i)_j.\tau,a_{ij})}
 {(D,p,B,U,K,\bot)\step{C-Constant}(D',p',B\cup\{k\},U,K[(i,j)\mapsto k],\bot)}

\Rule{C-Guard}{S\vdash i:\mathsf{ready}\quad c_i=\mathsf{if}\ e\quad
 (D_0,p_0)=\operatorname{unitize}(D,p)\quad\emit(D_0,\sigma_ep_0)=(D',p')}
 {S\step{C-Guard}(D',p',B,U\setminus\{i\},K,\bot)}

\Rule{C-Bind}{S\vdash i:\mathsf{ready}\quad c_i=\mathsf{let}\ p_0[:\tau]=e\\
 N=\binds(p_0)\quad(D',p')=\operatorname{typedBind}(D,p,p_0,\tau,e)}
 {S\step{C-Bind}(D',p',B\cup\set(N),U\setminus\{i\},K,\bot)}
\end{mathpar}
$\operatorname{typedBind}$ unitizes the prefix if necessary and emits Extend
with the literal Rust expression
\code{\{ let PAT: TYPE = EXPR; (BOUND_NAMES,) \}}; the annotation may be absent.
The written simple annotation supplies known output types; other pattern
bindings use local Rust inference. C-Constant uses this identical construction.

C-Match is a composed transition, not one opaque unexplained instruction.
Its argument context is $(D,p,\widehat B,L,M,E)$, with original prefix sort
$B_0$, earlier bare atom names $L$, positional mapping $M$, and pending local
equalities $E$. During this internal walk $\widehat B$ is the prospective
sort including not-yet-joined columns; it need not equal $\sort(p)$.
Only completed clause states assert that equality. The disjoint ordered cases
below define $\Sigma;B_0;K\vdash(a_j,j)\leadsto_A v_j$ and its context updates:
\begin{mathpar}

\Rule{A-Constant}{K(i,j)=k}{a_j\leadsto_A k\quad[\text{no expression emission}]}

\Rule{A-Variable}{a_j=x\quad x\notin L}{a_j\leadsto_A x\quad[L:=L\cup\{x\}]}

\Rule{A-Repeat}{a_j=x\in L\quad k\text{ fresh}}
 {a_j\leadsto_A k\quad[E:=E\concat(k=x)]}

\Rule{A-Wildcard}{a_j=\_\quad k\text{ fresh}}
 {a_j\leadsto_A k\quad[\text{retain its distinct witness column}]}

\Rule{A-Prefix}{a_j=e\quad\uses(e,L\setminus B_0)=\emptyset\quad k\text{ fresh}}
 {a_j\leadsto_A k\quad[p:=\operatorname{ext}_{\{k:\tau_j\},(e,)}(\operatorname{unitize}(p))]}

\Rule{A-Local}{a_j=e\quad\uses(e,L\setminus B_0)\ne\emptyset\quad k\text{ fresh}}
 {a_j\leadsto_A k\quad[E:=E\concat(k.\mathsf{eq}(\&e))]}
\end{mathpar}
A-Variable/Repeat require no K-entry; A-Wildcard/Prefix/Local additionally
exclude bare variables. Every case appends $a_j^{\rm schema}\mapsto v_j$ to
$M$ and adds $v_j$ to $\widehat B$. Prefix Extend retains old columns and
therefore creates a join key. The final C-Match rule is
\[
\Rule{C-Match}{S\vdash i:\mathsf{ready}\quad c_i=R(\bar a)\quad
 (D,p,B,\emptyset,\emptyset,\emptyseq)\leadsto_A^*(D_1,p_1,B_1,L,M,E)\\
 (D_2,r)=\emit(D_1,\delta_M(R))\quad
 (D_3,t)=\operatorname{joinPrefix}(D_2,p_1,r)\\
 (D',p')=\operatorname{selectList}(D_3,t,E)}
 {S\step{C-Match}(D',p',B_1,U\setminus\{i\},K,\bot)}.
\]
joinPrefix emits natural join when $p_1$ exists and otherwise returns $r$.
selectList emits the listed equalities in argument order. The diagnostic
trace records C-Match once; the A-judgment is its specified derivation, not a
second exported IR.

Define $\operatorname{keepHead}(B,h)$ as the union of head dependencies for
an entirely bare head, and as $B$ for a computed head. Then
\[
\Rule{C-Finish}{\operatorname{normalized}(K)\quad U=\emptyset\quad o=\bot\\
 W=\operatorname{keepHead}(B,h)
 \quad(D_0,p_0)=\operatorname{unitize}(D,p)\quad\emit(D_0,\pi_Wp_0)=(D',r)}
 {S\step{C-Finish}(D',r,W,\emptyset,K,\mathsf{output}\ r\ \mathsf{as}\ h(\bar e))}.
\]
Only a finished state can become an RA module. The implementation checks the
arity cap over the complete generated module before returning; unlike graph
emission it does not enforce that bound as a premise of each internal A-step.

\begin{theorem}[Prefix lowering, conditional set correctness]
Under $\mathcal H_{\rm rel}$ and CQ well-scoping, every complete prefix
derivation yields $\sem q_I$. Its policy terminates or rejects with a diagnostic.
\end{theorem}
\begin{proof}
For a written schedule, induct on complete clauses. Unit is the empty
environment. Literal normalization adds fixed typed coordinates; matching
against them enforces exactly the corresponding argument equalities. For an
atom, A-Variable supplies the row coordinate and joins any previous binding;
A-Repeat creates two positional coordinates followed by their equality;
A-Wildcard retains the input coordinate as an existential witness. For A-Prefix,
the expression depends only on the prefix environment, so extension followed
by equijoin enforces exactly its value constraint. For A-Local, the needed
earlier coordinate is available after Rename/Join and the selection enforces
the same constraint. The cases partition all arguments, proving C-Match.
Guard is selection. Bind extends by the exact pattern result and retains its
input environment. Thus the prefix, after forgetting generated witness names,
is the relational composition of executed clauses. Finish applies the head
function; retaining extra witnesses cannot change a set of pure head values.

For an alternative legal order, remove its first chosen clause from the
written pending list. It can be moved left across earlier pending clauses
because C-Ready requires all of them and itself to be movable. Movable atoms
contain only equational constraints, which commute under relational join.
A literal binding has no inputs and is moved only where its names remain
fresh; atoms using a preceding let's value have that name in $R_i$ and cannot
precede it. A moved guard has all inputs already provided, so it commutes with
intervening joins/bindings that add no required value. Under host admissibility
these exchanges have no effects or representation observations. Repeating
this adjacent-exchange argument gives the chosen order's set equality with
written order. The finite literal prefix terminates; each clause transition
removes one identity from $U$; Finish ends execution. For valid written scope
the smallest remaining clause is ready after normalization, so policy does
not get stuck. Replayed illegal traces may reject, as specified.
\end{proof}

Prefix policy selects only C-Ready candidates. Written minimizes source index;
other policies prioritize non-atom guards/bindings, then atoms sharing a key,
then (for GreedyWidth/Auto) the predicted union arity, then source index.
For atom $i$ the implemented estimate is
$|B|+a_i-|O_i\cap B|-|K_i|$, including anonymous/repeated positional columns.
Auto on this domain does not claim a join tree. The independent RA checker
reads emitted operations, but shares D-analysis; its independence does not
extend to proving the correctness of that analysis.
